\documentclass[10pt,a4paper]{article}
\usepackage[utf8]{inputenc}
\usepackage[T1]{fontenc}
\usepackage[brazil]{babel}
\usepackage{geometry}
\usepackage{hyperref}
\usepackage{booktabs}
\geometry{margin=1.65cm}
\hypersetup{colorlinks=true,urlcolor=blue}
\setlength{\parskip}{0.35em}
\setlength{\parindent}{0pt}

\begin{document}

\begin{center}
{\Large\textbf{GuiaMed --- Entrega de Andamento}}\\[0.25em]
{\small Projeto de Inteligência Artificial e Segurança e Otimização de BD}\\[0.45em]
\textbf{Grupo GuiaMed}\\
\small João Roberto Santos Cardoso --- 231070711; Gabriel de Souza Nascimento --- 231071404;\\
Daniel Queiroga Santana Martins --- 231070654; Pedro Henrique Freitas Couto --- 231071331;\\
Pedro Felipe dos Anjos Santos --- 232070649.\\[0.25em]
\url{https://github.com/Roberto-Cardoso-1/GuiaMed}
\end{center}

\section*{Parte comum --- relatório de andamento}

\textbf{O que funciona.} A V1 é um app Android nativo de acessibilidade para pessoas cegas ou com baixa visão. Um toque lê embalagem, dois leem receita, três registram alergia por voz e quatro repetem instruções. O retorno combina TTS, vibração e painel de alto contraste. CameraX e ML Kit fazem OCR no aparelho; o parser local consulta 10.777 denominações DCB/Anvisa, reconhece dose e validade (inclusive \texttt{V02 28}), lê receitas em várias linhas e alerta alergia. As funções principais e o armazenamento de alergias funcionam sem computador ou internet.

\textbf{Mudanças desde 27/10.} O protótipo virou uma V1 compilável em Kotlin, com fluxo contínuo opcional para virar a embalagem, compartilhamento de imagem/PDF, leitura melhorada de receitas, persistência local, API FastAPI opcional e testes automatizados. Para a entrega acadêmica foi executado um RAG local, n8n auto-hospedado e PostgreSQL com auditoria e menor privilégio.

\textbf{O que falta e riscos.} Falta ampliar o teste físico com mais aparelhos, embalagens reflexivas e receitas manuscritas. OCR ruim pode induzir identificação incorreta; por isso o app usa múltiplas capturas, busca conservadora e recusa quando a confiança é insuficiente. Informação de saúde nunca substitui médico ou farmacêutico.

\section*{Parte de IA}

O fluxo comprovado é \texttt{n8n local -> POST /rag/consultar -> FastAPI -> RAG TF-IDF -> resposta citada -> auditoria}. A base tem 24 bulas fictícias e 144 chunks (uma seção por chunk). O workflow versionado executou três nós com sucesso e retornou \texttt{status\_evidencia: PASSOU}, fonte \texttt{pressormin\#indicacao}, auditoria, latência e custo.

Os prompts v1 e v2 estão versionados. Cinco casos do RAG documentam entrada, esperado e obtido, com cinco aprovações. Foi executado fine-tuning local do \textbf{SmolLM2-135M-Instruct por LoRA}, com 252 exemplos de treino e 48 de teste (24 classes fictícias), três épocas e 189 atualizações. Na mesma avaliação de geração livre, o original acertou 7/48 (14,58\%) e o ajustado 30/48 (62,50\%), ganho de 47,92 pontos. A perda média caiu de 7,2616 para 0,4388. O treino durou 172,50 s em CPU. A avaliação por exemplo levou 866,97 ms no original e 383,84 ms no ajustado, incluindo perda e geração. O custo de API foi US\$ 0, sem medição elétrica. Dados, pesos LoRA, logs e todos os erros estão em \texttt{evidencias/ia/finetuning/}. Os templates sintéticos limitam a generalização. O RAG em memória manteve média abaixo de 1 ms em 100 consultas.

\newpage

\section*{Parte de Segurança e Otimização}

O PostgreSQL local usa os papéis \texttt{rl\_app\_ia}, \texttt{rl\_paciente}, \texttt{rl\_suporte} e \texttt{rl\_farmaceutico}. A IA só executa funções \texttt{SECURITY DEFINER} validadas e não recebe acesso ao schema privado. Uma tentativa de \texttt{SELECT} direto produziu \texttt{permission denied}; a função autorizada registrou alergia e retornou alerta verdadeiro. RLS mostrou 1 registro para o dispositivo A e 0 para B. A visão de suporte exibiu apenas o hash mascarado.

A carga real criou 50.000 chunks e 1.000.000 de eventos fictícios. Dois planos foram medidos com \texttt{EXPLAIN (ANALYZE, BUFFERS)}: Q1 caiu de 37,647 ms (\textit{Seq Scan}) para 0,120 ms com full-text search e GIN; Q2 caiu de 109,734 ms para 0,075 ms com índice parcial por data para alertas críticos.

\section*{Parte de integração}

O n8n consultou o RAG pela API e criou a auditoria \texttt{221bec78-044b-4f94-be66-c2ae8562ca36}. Pelo perfil de suporte, a prova mostra solicitante \texttt{********************************************************a0d9a3ee}, operação \texttt{rag\_local}, fontes consultadas, horário, latência de 0,444 ms e custo zero. O identificador original e a pergunta livre não são armazenados: somente SHA-256 e impressão digital. Assim, um perfil sem permissão não lê a tabela nem por meio da IA, mas a função mínima registra o necessário.

\section*{Contribuição individual --- João Roberto Santos Cardoso}

Implementei e integrei a V1 Android/FastAPI do GuiaMed, incluindo CameraX, OCR local, acessibilidade por TTS/vibração e contrato da API. Ampliei reconhecimento, receitas, validade e alergias, gerei o APK offline e os testes de regressão. Também implementei o RAG local, o workflow n8n, o experimento supervisionado e as evidências de segurança, otimização e auditoria.

\small\textit{Evidências completas e reproduzíveis em \texttt{evidencias/}; dados fictícios ou anonimizados; credenciais fora do repositório.}

\end{document}
